\subsection{A certified census over the primes $2,3,5,7$}

Put $S=\{2,3,5,7\}$. For an even integer $m$, represent an $S$-supported
class in $\mathbb Q^\times/\mathbb Q^{\times m}$ uniquely as
$\varepsilon2^{e_2}3^{e_3}5^{e_5}7^{e_7}$, where
$\varepsilon\in\{1,-1\}$ and $0\leq e_p<m$. For cubic classes the sign
is a cube, and we take $\varepsilon=1$. The tables below give every
coefficient class, including the positive-rank classes. A pair of classes
defines the punctured diagonal curve
\[
 U_{\alpha,\beta,m}:
 \alpha X^m+\beta Y^m=Z^m,\qquad XYZ\ne0,
 \qquad t=\alpha(X/Z)^m.
\]
Every statement about its section set in this subsection concerns full
profinite Selmer sections. The all-section extension requires A.20.

\begin{statement}{A.37}{A proof certificate for each computed analytic rank}{Proved}
The coefficient tables accompanying A.38, A.40 and A.41 determine the exact
analytic rank of every curve in them. Each nonzero derivative is enclosed
in a real ball that excludes zero, including a proved bound for the
infinite tail. Each asserted central zero follows from the exact root
number or an independently checked rational point of infinite order.
All exact coefficient vectors and the bounds used in the certificates are
preserved with the programs.
\end{statement}
\begin{proof}
We give the analytic and arithmetic certificates, so that the rank labels
do not depend on a numerical zero tolerance. PARI/GP supplies an integral
minimal model, its conductor $N$, root number $w$, and exact coefficients
$a_n$ of its Hasse--Weil $L$-function. These are the documented
\texttt{ellminimalmodel}, \texttt{ellglobalred}, \texttt{ellrootno}, and
\texttt{ellan} operations \citep{AMPARI2026}. The source archive used was
PARI/GP 2.19.1, with SHA256
\begin{center}\scriptsize\ttfamily
419553b11ff8c3ecb8e9031b7d91d7cee003a5d721e685fcb499565c96fe2943.
\end{center}
The software record and the complete build log are included. Arb, through
python-flint 0.9.0, supplies rigorous real balls and rigorous incomplete
gamma and exponential integral evaluations
\citep{AMJohansson2017,AMPythonFlint2026}.

Let $c=2\pi/\sqrt N$ and $r=e^{-c}$. Hasse's bound and the Euler-factor
recurrences give
\[
 |a_n|\leq d(n)\sqrt n\leq2n.
\]
Here the divisor bound follows by pairing divisors below and above
$\sqrt n$. Modularity and the functional equation give, by splitting
the Mellin integral at $1/\sqrt N$,
\begin{align*}
 w=1:\quad L(E,1)&=2\sum_{n\geq1}\frac{a_n}{n}r^n,\\
 w=-1:\quad L(E,1)&=0,\qquad
 L'(E,1)=2\sum_{n\geq1}\frac{a_n}{n}E_1(cn),
 \quad E_1(x)=\int_x^\infty e^{-u}\frac{du}{u}.
\end{align*}
The absolute tails after $M$ terms are at most, respectively,
\[
 T_0=\frac{4r^{M+1}}{1-r},\qquad
 T_1=\frac{4r^{M+1}}{c(M+1)(1-r)}.
\]
The second estimate uses $E_1(x)\leq e^{-x}/x$. We initially take
$M=8(\lfloor\sqrt N\rfloor+1)$ and 160-bit balls. Adding the symmetric
ball of radius an upper bound for $T_i$ gives a rigorous enclosure of the
complete value. A ball excluding zero directly proves nonvanishing;
no inference from a small floating-point value or an unproved denominator
bound is made.

If $w=1$ and this interval contains zero, the program obtains a rational
point $P$ and verifies its equation using exact rational arithmetic. At
the good inert primes $11,19$, every quartic-coefficient curve has
$12,20$ points, so its rational torsion has order dividing four. At
$11,17$, every cubic or sextic-coefficient curve has $12,18$ points, so
its rational torsion has order dividing six. These reductions are good
because the discriminants are supported on $S$. The assertion on torsion
uses injectivity on torsion of order prime to the reduction characteristic
at each of the two different primes; the two orders exclude torsion at
the exceptional characteristics as well. We check $4P\ne O$ or $6P\ne O$
by the exact addition law. Thus $P$ has infinite order.
For these CM curves, \citet[Corollary C and Remark (3), pp.527--528]{AMRubin1987}
then implies $L(E,1)=0$: a nonzero central value would imply algebraic
rank zero. The even functional equation implies order at least two.

For the higher derivative put
\[
 \Lambda(s)=c^{-s}\Gamma(s)L(E,s),\qquad h=2^{-45}.
\]
The split Mellin integral gives the exact identity
\[
 \Lambda(1+h)=\sum_{n\geq1}a_n\left\{
 \frac{\Gamma(1+h,cn)}{(cn)^{1+h}}+
 w\frac{\Gamma(1-h,cn)}{(cn)^{1-h}}\right\}.
\]
We evaluate through $M=20(\lfloor\sqrt N\rfloor+1)$ at 256 bits. For
$|h|\leq1$, the absolute tail is at most
\[
 T=\frac{4r^{M+1}}{1-r}
       \left(\frac1c+\frac1{c^2(M+1)}\right),
\]
since the two powers in the integrand are bounded by $2t$, for $t\geq1$.
For every integer $j\geq1$ and $|x|\leq h\leq1$ we also have
\[
 |\Lambda^{(j)}(1+x)|\leq B_j,
 \qquad
 B_j=\frac{4j^j r}{(1-r)^2}
       \left(\frac1c+\frac2{c^2}+\frac2{c^3}\right).
\]
Indeed $|f(it/\sqrt N)|\leq2e^{-ct}/(1-e^{-ct})^2$,
$t^{|x|}\leq t$, and $(\log t)^j\leq j^j t$ for $t\geq1$.
Integrating $t^2e^{-ct}$ proves the displayed bound.

If the known vanishing and parity give order at least $k$, Taylor's
theorem therefore encloses the first possible derivative by
\[
 L^{(k)}(E,1)\ \in\
 \frac{c k!\Lambda(1+h)}{h^k}
 +\left[-\frac{cB_{k+2}h^2}{(k+1)(k+2)},
         \frac{cB_{k+2}h^2}{(k+1)(k+2)}\right].
\]
Here the lower derivatives vanish, the derivative of order $k+1$ vanishes
by the functional equation, and $\Lambda^{(k)}(1)=L^{(k)}(E,1)/c$.
Every case needing this step has $k=2$ and an interval excluding zero.
Thus all the analytic ranks in this census are exactly $0,1$, or $2$.

The files retain each minimal model and change of variables, every exact
$a_n$ used, the torsion calculation, each rational point certificate,
the ball precision, truncation length, all tail bounds, and the final
interval. The supplied code rechecks the rational points without using
PARI's numerical analytic-rank routine. This proves the claimed type of
certificate for every row.
\end{proof}

\begin{statement}{A.38}{All 81 cubic coefficient classes}{Proved}
For $d=2^{e_2}3^{e_3}5^{e_5}7^{e_7}$, $0\leq e_p<3$, let
$E_d:y^2=x^3-432d^2$. Of the 81 curves, 34 have analytic rank zero,
41 have analytic rank one, and six have analytic rank two. The rank-two
values of $d$ are
\[
 30,\quad126,\quad210,\quad420,\quad2100,\quad2205.
\]
The following table gives every class. The CSV also gives each exact
conductor, root number, minimal model, and nonvanishing certificate.
\end{statement}
\begin{proof}
The program \texttt{a\_main\_cubic\_census.py} enumerates the $3^4$
exponent vectors, and \texttt{a\_main\_cubic\_rank2.py} supplies the six
second-derivative certificates by A.37. Every other case is certified by
the central value or first derivative. The exact output is
\texttt{a\_main\_cubic\_census.csv}; the separate rank-two CSV records
the completed value, tail, fourth-derivative bound and resulting interval.
The saved coefficient vectors suffice to repeat every enclosure without
recomputing any Fourier coefficient.
\end{proof}
{
\scriptsize
\setlength{\tabcolsep}{2pt}
\begin{longtable}{rrr|rrr|rrr|rrr}
\caption{Complete certified cubic analytic ranks; rows give $(\varepsilon,e_2,e_3)$ and columns give $(e_5,e_7)$.}\\
\toprule
$\varepsilon$&$e_2$&$e_3$&\multicolumn{3}{c}{$e_5=0$}&\multicolumn{3}{c}{$e_5=1$}&\multicolumn{3}{c}{$e_5=2$}\\
&&&0&1&2&0&1&2&0&1&2\\
\midrule
\endfirsthead
\toprule
$\varepsilon$&$e_2$&$e_3$&\multicolumn{3}{c}{$e_5=0$}&\multicolumn{3}{c}{$e_5=1$}&\multicolumn{3}{c}{$e_5=2$}\\
&&&0&1&2&0&1&2&0&1&2\\
\midrule
\endhead
\bottomrule\endfoot
+1&0&0&0&1&1&0&1&0&0&0&1\\
+1&0&1&0&0&0&1&1&1&1&1&1\\
+1&0&2&1&1&1&0&0&2&0&0&0\\
+1&1&0&0&0&1&0&1&1&1&0&1\\
+1&1&1&1&1&1&2&2&0&0&0&0\\
+1&1&2&0&2&0&1&1&1&1&1&1\\
+1&2&0&0&1&0&1&1&0&0&1&1\\
+1&2&1&1&1&1&0&2&0&0&2&0\\
+1&2&2&0&0&0&1&1&1&1&1&1\\
\end{longtable}
}


\begin{statement}{A.39}{The complete section lists in the cubic rank-zero classes}{Proved}
Among the $81^2=6561$ ordered pairs $(\alpha,\beta)$, the cubic
rank-zero criterion applies to 2754 pairs. Exactly 2654 have no Selmer
sections, 97 have only rational cusp packets, and three have one ordinary
point section and one cusp packet. The latter pairs, ordinary points and
tripod values are
\[
\begin{array}{c|c|c}
(\alpha,\beta)&[X:Y:Z]&t\\\hline
(1,2)&[-1:1:1]&-1\\
(2,1)&[1:-1:1]&2\\
(4,4)&[1:1:2]&1/2.
\end{array}
\]
The complete 6561-row CSV gives the criterion, every ordinary point, every
rational cusp, and the Kummer congruence defining the image of each cusp
packet on the tripod.
\end{statement}
\begin{proof}
Put $d=[\alpha\beta]_3$, using the normalized cube-free representative.
By A.27 the Jacobian is $E_d$. The torsion bound in A.37 and the division
polynomials give
\[
 E_d(\mathbb Q)_{\rm tors}=\begin{cases}
 \mathbb Z/3,&d=1,\\
 \mathbb Z/2,&d=2,\\
 0,&d\ne1,2.
 \end{cases}
\]
For completeness, rational two-torsion requires $x^3=432d^2$, which
holds exactly when $d=2$. The three-division polynomial is
$3x(x^3-1728d^2)$. The root $x=0$ has negative right side in the
curve equation; the other factor has a rational root exactly when $d=1$,
and then gives rational three-torsion. The reductions already bound the
whole torsion order by six.

Suppose the rank is zero and $d\ne1,2$. If $P$ is a rational point of
the diagonal cubic and $H$ its plane hyperplane class, then
$3[P]-H=0$ in its trivial rational Jacobian. Thus $P$ is a flex. More
explicitly, the complete degree-three plane linear system contains the
divisor $3P$, so a line cuts out $3P$. The Hessian of
$\alpha X^3+\beta Y^3-Z^3$ is
$-216\alpha\beta XYZ$. Its flexes are exactly the coordinate cusps.

If $d=1$ and $(\alpha,\beta)\ne(1,1)$, at some prime the three
coefficient valuations are distinct modulo three. In a putative local
point the three summands then have distinct valuations modulo three;
the least valuation occurs only once, an impossibility. For $(1,1)$,
the rational cusp gives an origin and the three torsion points are
exactly the three cusps. If $d=2$, the same obstruction excludes every
pair except $(1,2),(2,1),(4,4)$. At primes other than two it forces both
exponents to vanish; at two the allowed exponent pairs are
$(0,1),(1,0),(2,2)$. Each remaining curve has the displayed ordinary
point and a rational cusp, and its rational group has order two. These
are all its rational points.

At $t=0$ a rational cusp exists precisely when $\beta=1$; at $t=1$
precisely when $\alpha=1$; at infinity precisely when $\alpha=\beta$.
They are respectively $[0:1:1]$, $[1:0:1]$, and $[1:-1:0]$.
The packet images satisfy
\[
 [t]_3=[\alpha]_3,\qquad
 [1-t]_3=[\beta]_3,\qquad
 [1/t]_3=[\alpha^{-1}]_3,
\]
respectively. Each rational cusp upstairs contributes its whole packet
$H^1(\mathbb Q,\widehat{\mathbb Z}(1))$ after choosing a tangent origin.
A.24 converts the rational point list into the complete Selmer section
list. The program \texttt{a\_main\_cubic\_pairs.py} verifies the Hessian,
torsion equations, local obstruction, and every listed point, then
enumerates all ordered pairs. Each of the 34 rank-zero values of $d$
has 81 inverse pairs, giving $34\cdot81=2754$. The stated disjoint
counts follow from the three cusp conditions and the three exceptional
ordinary pairs. A blank classification outside these 2754 pairs makes
no assertion of emptiness.
\end{proof}

\begin{statement}{A.40}{All 512 quartic coefficient classes}{Proved}
For $A=\varepsilon2^{e_2}3^{e_3}5^{e_5}7^{e_7}$, $0\leq e_p<4$,
the curves $y^2=x^3+Ax$ have the following exact analytic-rank counts:
220 of rank zero, 256 of rank one, and 36 of rank two. Every coefficient
class is included in the following table and the corresponding CSV.
\end{statement}
\begin{proof}
The enumeration in \texttt{a\_main\_cm\_census.py quartic} has exactly
$2\cdot4^4=512$ inputs. It applies A.37 to each input, and records a
nonzero certified derivative of the asserted order. The 36 even-sign
central zeros have exact non-torsion point certificates and nonzero
second-derivative intervals. The file
\texttt{a\_main\_quartic\_census.csv} and its raw data directory preserve
all these certificates.
\end{proof}
{
\scriptsize
\setlength{\tabcolsep}{2pt}
\begin{longtable}{rrr|rrrr|rrrr|rrrr|rrrr}
\caption{Complete certified quartic analytic ranks; rows give $(\varepsilon,e_2,e_3)$ and columns give $(e_5,e_7)$.}\\
\toprule
$\varepsilon$&$e_2$&$e_3$&\multicolumn{4}{c}{$e_5=0$}&\multicolumn{4}{c}{$e_5=1$}&\multicolumn{4}{c}{$e_5=2$}&\multicolumn{4}{c}{$e_5=3$}\\
&&&0&1&2&3&0&1&2&3&0&1&2&3&0&1&2&3\\
\midrule
\endfirsthead
\toprule
$\varepsilon$&$e_2$&$e_3$&\multicolumn{4}{c}{$e_5=0$}&\multicolumn{4}{c}{$e_5=1$}&\multicolumn{4}{c}{$e_5=2$}&\multicolumn{4}{c}{$e_5=3$}\\
&&&0&1&2&3&0&1&2&3&0&1&2&3&0&1&2&3\\
\midrule
\endhead
\bottomrule\endfoot
+1&0&0&0&0&1&0&1&1&0&1&0&1&1&1&1&0&0&0\\
+1&0&1&1&1&0&1&1&2&2&0&0&1&1&1&0&0&1&2\\
+1&0&2&1&2&0&0&0&1&1&1&1&1&2&1&0&0&1&0\\
+1&0&3&0&1&1&1&0&2&1&0&1&1&0&1&1&0&0&0\\
+1&1&0&0&2&1&0&0&0&1&0&0&0&1&0&0&0&1&0\\
+1&1&1&0&0&1&0&0&0&1&0&2&0&1&0&0&0&1&0\\
+1&1&2&1&1&0&1&1&1&0&1&1&1&0&1&1&1&0&1\\
+1&1&3&0&0&1&0&0&0&1&0&0&0&1&0&0&0&1&0\\
+1&2&0&0&1&1&1&1&0&0&0&1&1&0&1&0&0&1&0\\
+1&2&1&0&1&1&1&1&1&0&1&0&2&1&0&1&0&0&0\\
+1&2&2&0&0&1&0&1&1&2&1&1&0&0&0&0&1&1&1\\
+1&2&3&0&0&1&0&1&0&0&0&0&1&1&1&1&1&0&1\\
+1&3&0&1&1&0&1&1&1&0&1&1&1&2&1&1&1&0&1\\
+1&3&1&1&1&0&1&1&1&2&1&1&1&0&1&1&1&2&1\\
+1&3&2&0&0&1&0&2&2&1&0&0&2&1&0&0&2&1&0\\
+1&3&3&1&1&0&1&1&1&0&1&1&1&2&1&1&1&0&1\\
-1&0&0&0&1&1&1&1&0&0&0&1&1&0&1&0&0&1&0\\
-1&0&1&0&1&1&1&1&1&0&1&0&2&1&0&1&0&0&0\\
-1&0&2&0&0&1&0&1&1&2&1&1&0&0&0&0&1&1&1\\
-1&0&3&0&0&1&0&1&0&0&0&0&1&1&1&1&1&0&1\\
-1&1&0&1&1&0&1&1&1&0&1&1&1&2&1&1&1&0&1\\
-1&1&1&1&1&0&1&1&1&2&1&1&1&0&1&1&1&2&1\\
-1&1&2&0&0&1&0&2&2&1&0&0&2&1&0&0&2&1&0\\
-1&1&3&1&1&0&1&1&1&0&1&1&1&2&1&1&1&0&1\\
-1&2&0&0&0&1&0&1&1&0&1&0&1&1&1&1&0&0&0\\
-1&2&1&1&1&0&1&1&2&2&0&0&1&1&1&0&0&1&2\\
-1&2&2&1&2&0&0&0&1&1&1&1&1&2&1&0&0&1&0\\
-1&2&3&0&1&1&1&0&2&1&0&1&1&0&1&1&0&0&0\\
-1&3&0&0&2&1&0&0&0&1&0&0&0&1&0&0&0&1&0\\
-1&3&1&0&0&1&0&0&0&1&0&2&0&1&0&0&0&1&0\\
-1&3&2&1&1&0&1&1&1&0&1&1&1&0&1&1&1&0&1\\
-1&3&3&0&0&1&0&0&0&1&0&0&0&1&0&0&0&1&0\\
\end{longtable}
}


\begin{statement}{A.41}{All 2592 sextic coefficient classes}{Proved}
For $B=\varepsilon2^{e_2}3^{e_3}5^{e_5}7^{e_7}$, $0\leq e_p<6$,
the curves $y^2=x^3+B$ have the following exact analytic-rank counts:
1136 of rank zero, 1296 of rank one, and 160 of rank two. Every
coefficient class is included in the following table and the corresponding
CSV. There are no unresolved rank labels in this census.
\end{statement}
\begin{proof}
The run \texttt{a\_main\_cm\_census.py sextic} enumerates all
$2\cdot6^4=2592$ inputs and applies A.37. The 160 even-sign central
zeros have exact non-torsion points and certified nonzero second
derivatives; all other rows have a certified nonzero value or first
derivative. See \texttt{a\_main\_sextic\_census.csv} and the raw data
directory. The same exact coefficient curves supply all ten target
families of A.34 after reduction modulo rational sixth powers.
\end{proof}
{
\scriptsize
\setlength{\tabcolsep}{2pt}
\begin{longtable}{rrr|rrrrrr|rrrrrr|rrrrrr|rrrrrr|rrrrrr|rrrrrr}
\caption{Complete certified sextic analytic ranks; rows give $(\varepsilon,e_2,e_3)$ and columns give $(e_5,e_7)$.}\\
\toprule
$\varepsilon$&$e_2$&$e_3$&\multicolumn{6}{c}{$e_5=0$}&\multicolumn{6}{c}{$e_5=1$}&\multicolumn{6}{c}{$e_5=2$}&\multicolumn{6}{c}{$e_5=3$}&\multicolumn{6}{c}{$e_5=4$}&\multicolumn{6}{c}{$e_5=5$}\\
&&&0&1&2&3&4&5&0&1&2&3&4&5&0&1&2&3&4&5&0&1&2&3&4&5&0&1&2&3&4&5&0&1&2&3&4&5\\
\midrule
\endfirsthead
\toprule
$\varepsilon$&$e_2$&$e_3$&\multicolumn{6}{c}{$e_5=0$}&\multicolumn{6}{c}{$e_5=1$}&\multicolumn{6}{c}{$e_5=2$}&\multicolumn{6}{c}{$e_5=3$}&\multicolumn{6}{c}{$e_5=4$}&\multicolumn{6}{c}{$e_5=5$}\\
&&&0&1&2&3&4&5&0&1&2&3&4&5&0&1&2&3&4&5&0&1&2&3&4&5&0&1&2&3&4&5&0&1&2&3&4&5\\
\midrule
\endhead
\bottomrule\endfoot
+1&0&0&0&0&0&1&1&1&1&1&0&0&0&1&0&0&1&1&1&0&0&1&1&1&0&0&1&0&0&0&1&1&0&0&0&1&1&1\\
+1&0&1&1&0&1&0&1&0&2&1&0&1&2&1&0&1&0&1&0&1&0&1&0&1&0&1&0&1&0&1&0&1&0&1&0&1&0&1\\
+1&0&2&1&2&1&0&1&0&0&1&0&1&0&1&2&1&2&1&0&1&0&1&0&1&0&1&0&1&0&1&0&1&0&1&0&1&0&1\\
+1&0&3&0&1&1&1&0&0&0&1&1&1&0&0&1&1&1&0&0&0&0&2&2&1&1&1&0&0&1&1&1&0&1&1&0&0&0&1\\
+1&0&4&0&1&2&1&0&1&1&0&1&0&1&0&1&0&1&0&1&0&1&0&1&0&1&0&1&2&1&0&1&0&1&0&1&0&1&0\\
+1&0&5&0&1&0&1&0&1&1&0&1&0&1&0&1&0&1&0&1&0&1&0&1&0&1&0&1&0&1&0&1&0&1&0&1&0&1&0\\
+1&1&0&1&0&1&1&0&1&1&0&1&1&0&1&1&2&0&1&2&0&0&1&1&0&1&1&0&0&1&0&0&1&1&1&0&1&1&0\\
+1&1&1&0&0&2&2&0&0&1&1&1&1&1&1&1&1&1&1&1&1&1&1&1&1&1&1&1&1&1&1&1&1&1&1&1&1&1&1\\
+1&1&2&1&1&1&1&1&1&0&0&0&2&0&0&0&2&0&0&0&0&2&2&0&2&2&0&0&0&0&0&0&0&0&0&0&0&2&0\\
+1&1&3&1&0&0&1&0&0&0&0&1&0&0&1&1&1&0&1&1&0&0&1&2&0&1&0&1&0&1&1&0&1&1&0&0&1&0&0\\
+1&1&4&0&0&2&0&0&0&1&1&1&1&1&1&1&1&1&1&1&1&1&1&1&1&1&1&1&1&1&1&1&1&1&1&1&1&1&1\\
+1&1&5&1&1&1&1&1&1&0&0&0&0&0&0&0&2&0&0&0&0&0&0&0&0&0&0&0&0&0&2&0&0&0&0&0&0&2&0\\
+1&2&0&0&1&1&1&0&0&0&0&0&1&1&1&1&1&1&0&0&0&1&1&0&0&0&1&0&0&1&1&1&0&0&1&1&1&0&0\\
+1&2&1&1&0&1&0&1&0&0&1&2&1&0&1&0&1&0&1&0&1&2&1&2&1&0&1&0&1&0&1&0&1&0&1&0&1&0&1\\
+1&2&2&1&2&1&2&1&0&0&1&0&1&0&1&0&1&2&1&0&1&0&1&0&1&0&1&0&1&2&1&0&1&0&1&0&1&0&1\\
+1&2&3&1&1&0&0&0&1&1&1&2&2&0&1&1&0&0&0&1&1&0&1&1&1&2&0&1&1&1&0&0&0&0&0&0&1&1&1\\
+1&2&4&0&1&0&1&0&1&1&0&1&0&1&0&1&0&1&2&1&0&1&0&1&0&1&0&1&0&1&0&1&2&1&0&1&0&1&0\\
+1&2&5&0&1&0&1&0&1&1&0&1&0&1&0&1&0&1&0&1&0&1&0&1&0&1&0&1&0&1&0&1&0&1&0&1&0&1&0\\
+1&3&0&1&1&2&1&1&0&1&1&0&1&1&0&0&1&2&2&1&0&1&0&1&1&0&1&1&2&0&1&0&0&0&1&1&0&1&1\\
+1&3&1&2&2&0&0&0&0&1&1&1&1&1&1&1&1&1&1&1&1&1&1&1&1&1&1&1&1&1&1&1&1&1&1&1&1&1&1\\
+1&3&2&1&1&1&1&1&1&2&2&2&0&0&0&0&0&0&2&2&0&2&2&0&0&0&2&0&2&2&0&0&0&0&0&0&0&0&0\\
+1&3&3&0&1&0&0&1&0&1&0&2&1&2&0&2&1&1&0&1&1&0&0&1&0&0&1&1&1&2&1&1&0&0&1&0&0&1&0\\
+1&3&4&0&0&0&0&2&0&1&1&1&1&1&1&1&1&1&1&1&1&1&1&1&1&1&1&1&1&1&1&1&1&1&1&1&1&1&1\\
+1&3&5&1&1&1&1&1&1&0&0&0&0&0&0&0&2&0&2&0&0&0&0&2&0&2&0&0&0&0&2&0&2&0&0&0&0&0&0\\
+1&4&0&0&1&1&0&1&1&1&1&0&1&1&0&0&0&1&0&0&1&1&0&1&1&0&1&0&1&0&0&1&0&0&1&1&0&1&1\\
+1&4&1&1&1&1&1&1&1&0&2&0&0&2&0&0&2&0&0&0&0&0&0&0&2&0&0&0&2&0&0&0&0&0&0&0&0&0&0\\
+1&4&2&0&0&0&0&0&0&1&1&1&1&1&1&1&1&1&1&1&1&1&1&1&1&1&1&1&1&1&1&1&1&1&1&1&1&1&1\\
+1&4&3&0&1&0&0&1&0&0&1&0&0&1&0&0&1&1&2&1&1&1&0&0&1&0&0&1&1&0&1&1&0&0&0&1&0&0&1\\
+1&4&4&1&1&1&1&1&1&0&0&0&0&0&0&0&2&0&0&2&0&0&0&0&0&0&0&0&0&0&0&0&0&0&0&0&0&0&0\\
+1&4&5&0&0&0&0&0&0&1&1&1&1&1&1&1&1&1&1&1&1&1&1&1&1&1&1&1&1&1&1&1&1&1&1&1&1&1&1\\
+1&5&0&0&1&1&0&1&1&0&1&1&0&1&1&0&0&1&0&0&1&1&1&0&1&1&0&0&1&0&2&1&0&1&0&1&1&0&1\\
+1&5&1&0&0&0&0&0&0&1&1&1&1&1&1&1&1&1&1&1&1&1&1&1&1&1&1&1&1&1&1&1&1&1&1&1&1&1&1\\
+1&5&2&1&1&1&1&1&1&0&0&2&0&0&2&0&0&2&0&0&2&0&0&0&2&0&0&0&0&0&0&0&0&0&0&2&0&0&0\\
+1&5&3&0&0&1&0&0&1&0&1&0&0&1&0&1&0&1&1&0&1&1&0&0&1&0&0&0&1&1&0&1&1&0&0&1&0&0&1\\
+1&5&4&0&0&0&0&0&0&1&1&1&1&1&1&1&1&1&1&1&1&1&1&1&1&1&1&1&1&1&1&1&1&1&1&1&1&1&1\\
+1&5&5&1&1&1&1&1&1&0&0&2&0&0&0&0&0&0&0&0&0&0&0&0&0&0&0&0&0&0&0&0&0&0&0&0&0&0&0\\
-1&0&0&0&1&1&1&0&0&0&1&1&1&0&0&1&1&1&0&0&0&0&2&2&1&1&1&0&0&1&1&1&0&1&1&0&0&0&1\\
-1&0&1&0&1&2&1&0&1&1&0&1&0&1&0&1&0&1&0&1&0&1&0&1&0&1&0&1&2&1&0&1&0&1&0&1&0&1&0\\
-1&0&2&0&1&0&1&0&1&1&0&1&0&1&0&1&0&1&0&1&0&1&0&1&0&1&0&1&0&1&0&1&0&1&0&1&0&1&0\\
-1&0&3&0&0&0&1&1&1&1&1&0&0&0&1&0&0&1&1&1&0&0&1&1&1&0&0&1&0&0&0&1&1&0&0&0&1&1&1\\
-1&0&4&1&0&1&0&1&0&2&1&0&1&2&1&0&1&0&1&0&1&0&1&0&1&0&1&0&1&0&1&0&1&0&1&0&1&0&1\\
-1&0&5&1&2&1&0&1&0&0&1&0&1&0&1&2&1&2&1&0&1&0&1&0&1&0&1&0&1&0&1&0&1&0&1&0&1&0&1\\
-1&1&0&1&0&0&1&0&0&0&0&1&0&0&1&1&1&0&1&1&0&0&1&2&0&1&0&1&0&1&1&0&1&1&0&0&1&0&0\\
-1&1&1&0&0&2&0&0&0&1&1&1&1&1&1&1&1&1&1&1&1&1&1&1&1&1&1&1&1&1&1&1&1&1&1&1&1&1&1\\
-1&1&2&1&1&1&1&1&1&0&0&0&0&0&0&0&2&0&0&0&0&0&0&0&0&0&0&0&0&0&2&0&0&0&0&0&0&2&0\\
-1&1&3&1&0&1&1&0&1&1&0&1&1&0&1&1&2&0&1&2&0&0&1&1&0&1&1&0&0&1&0&0&1&1&1&0&1&1&0\\
-1&1&4&0&0&2&2&0&0&1&1&1&1&1&1&1&1&1&1&1&1&1&1&1&1&1&1&1&1&1&1&1&1&1&1&1&1&1&1\\
-1&1&5&1&1&1&1&1&1&0&0&0&2&0&0&0&2&0&0&0&0&2&2&0&2&2&0&0&0&0&0&0&0&0&0&0&0&2&0\\
-1&2&0&1&1&0&0&0&1&1&1&2&2&0&1&1&0&0&0&1&1&0&1&1&1&2&0&1&1&1&0&0&0&0&0&0&1&1&1\\
-1&2&1&0&1&0&1&0&1&1&0&1&0&1&0&1&0&1&2&1&0&1&0&1&0&1&0&1&0&1&0&1&2&1&0&1&0&1&0\\
-1&2&2&0&1&0&1&0&1&1&0&1&0&1&0&1&0&1&0&1&0&1&0&1&0&1&0&1&0&1&0&1&0&1&0&1&0&1&0\\
-1&2&3&0&1&1&1&0&0&0&0&0&1&1&1&1&1&1&0&0&0&1&1&0&0&0&1&0&0&1&1&1&0&0&1&1&1&0&0\\
-1&2&4&1&0&1&0&1&0&0&1&2&1&0&1&0&1&0&1&0&1&2&1&2&1&0&1&0&1&0&1&0&1&0&1&0&1&0&1\\
-1&2&5&1&2&1&2&1&0&0&1&0&1&0&1&0&1&2&1&0&1&0&1&0&1&0&1&0&1&2&1&0&1&0&1&0&1&0&1\\
-1&3&0&0&1&0&0&1&0&1&0&2&1&2&0&2&1&1&0&1&1&0&0&1&0&0&1&1&1&2&1&1&0&0&1&0&0&1&0\\
-1&3&1&0&0&0&0&2&0&1&1&1&1&1&1&1&1&1&1&1&1&1&1&1&1&1&1&1&1&1&1&1&1&1&1&1&1&1&1\\
-1&3&2&1&1&1&1&1&1&0&0&0&0&0&0&0&2&0&2&0&0&0&0&2&0&2&0&0&0&0&2&0&2&0&0&0&0&0&0\\
-1&3&3&1&1&2&1&1&0&1&1&0&1&1&0&0&1&2&2&1&0&1&0&1&1&0&1&1&2&0&1&0&0&0&1&1&0&1&1\\
-1&3&4&2&2&0&0&0&0&1&1&1&1&1&1&1&1&1&1&1&1&1&1&1&1&1&1&1&1&1&1&1&1&1&1&1&1&1&1\\
-1&3&5&1&1&1&1&1&1&2&2&2&0&0&0&0&0&0&2&2&0&2&2&0&0&0&2&0&2&2&0&0&0&0&0&0&0&0&0\\
-1&4&0&0&1&0&0&1&0&0&1&0&0&1&0&0&1&1&2&1&1&1&0&0&1&0&0&1&1&0&1&1&0&0&0&1&0&0&1\\
-1&4&1&1&1&1&1&1&1&0&0&0&0&0&0&0&2&0&0&2&0&0&0&0&0&0&0&0&0&0&0&0&0&0&0&0&0&0&0\\
-1&4&2&0&0&0&0&0&0&1&1&1&1&1&1&1&1&1&1&1&1&1&1&1&1&1&1&1&1&1&1&1&1&1&1&1&1&1&1\\
-1&4&3&0&1&1&0&1&1&1&1&0&1&1&0&0&0&1&0&0&1&1&0&1&1&0&1&0&1&0&0&1&0&0&1&1&0&1&1\\
-1&4&4&1&1&1&1&1&1&0&2&0&0&2&0&0&2&0&0&0&0&0&0&0&2&0&0&0&2&0&0&0&0&0&0&0&0&0&0\\
-1&4&5&0&0&0&0&0&0&1&1&1&1&1&1&1&1&1&1&1&1&1&1&1&1&1&1&1&1&1&1&1&1&1&1&1&1&1&1\\
-1&5&0&0&0&1&0&0&1&0&1&0&0&1&0&1&0&1&1&0&1&1&0&0&1&0&0&0&1&1&0&1&1&0&0&1&0&0&1\\
-1&5&1&0&0&0&0&0&0&1&1&1&1&1&1&1&1&1&1&1&1&1&1&1&1&1&1&1&1&1&1&1&1&1&1&1&1&1&1\\
-1&5&2&1&1&1&1&1&1&0&0&2&0&0&0&0&0&0&0&0&0&0&0&0&0&0&0&0&0&0&0&0&0&0&0&0&0&0&0\\
-1&5&3&0&1&1&0&1&1&0&1&1&0&1&1&0&0&1&0&0&1&1&1&0&1&1&0&0&1&0&2&1&0&1&0&1&1&0&1\\
-1&5&4&0&0&0&0&0&0&1&1&1&1&1&1&1&1&1&1&1&1&1&1&1&1&1&1&1&1&1&1&1&1&1&1&1&1&1&1\\
-1&5&5&1&1&1&1&1&1&0&0&2&0&0&2&0&0&2&0&0&2&0&0&0&2&0&0&0&0&0&0&0&0&0&0&2&0&0&0\\
\end{longtable}
}


\begin{statement}{A.42}{Full quartic and sextic pair tables and their section sets}{Proved}
The supplied full pair CSVs have 262144 quartic rows and 6718464 sextic
rows, compressed losslessly by gzip. The disjoint classification counts
are
\[
\begin{array}{l|r|r}
 &m=4&m=6\\\hline
\text{all ordered pairs}&262144&6718464\\
\text{at least one rank-zero target}&213496&6689532\\
\text{no Selmer sections}&212200&6681792\\
\text{cusp packets only}&1293&7725\\
\text{ordinary point sections only}&3&15\\
\text{ordinary points and cusp packets}&0&0\\
\text{criterion does not settle the pair}&48648&28932.
\end{array}
\]
Every ordinary point and every cusp packet in a settled row is specified
exactly by the accompanying dictionaries and the decoding rule below.
\end{statement}
\begin{proof}
The coefficient of each target is the indicated element of
$\mathbb Q^\times/\mathbb Q^{\times m}$. In the order used by the mask
bits, the quartic targets of A.25 are
\[
 -4\alpha\beta,\qquad
 4\alpha/\beta^2,\qquad
 4\beta/\alpha^2.
\]
The sextic targets of A.34, with bit numbers $0$ through $9$, are
\[
\begin{array}{c|l@{\qquad}c|l}
0&\alpha^2\beta^3&1&\alpha^3\beta^2\\
2&\alpha^2\beta&3&\alpha\beta^2\\
4&-\alpha^3\beta&5&-\alpha\beta^3\\
6&-4\alpha\beta&7&4\alpha/\beta^2\\
8&4\beta/\alpha^2&9&-432(\alpha\beta)^2.
\end{array}
\]
A bit equals one exactly when that target has analytic rank zero in
A.40 or A.41. Thus a nonzero mask is exactly the union of the stated
criteria, with no omitted target or coefficient class.

The class dictionary uses the integer identifier
\[
 \operatorname{id}(\varepsilon,e_2,e_3,e_5,e_7)
 =\frac{1-\varepsilon}{2}m^4+e_2m^3+e_3m^2+e_5m+e_7.
\]
The full CSV columns are
\texttt{alpha\_id,beta\_id,criterion\_mask,point\_ids,cusp\_flags}.
To recover a target class one adds the coefficient exponent vectors
modulo $m$ and the sign exponents modulo two. This is a finite exact join
between the full coefficient-rank tables printed above and the displayed
target rules. It is also a compact presentation in the paper of every
one of the millions of CSV rows.

For a nonzero mask, A.35 restricts ordinary tripod values to
\[
\begin{split}
m=4:\quad&\{-1,1/2,2\},\\
m=6:\quad&\{-8,-3,-1,-1/3,-1/8,1/9,1/4,1/2,3/4,8/9,
                 9/8,4/3,2,4,9\}.
\end{split}
\]
For each such value $t$, form its class pair
$([t]_m,[1-t]_m)$. It occurs in precisely that normalized pair of
classes. If $\alpha,\beta$ are its selected representatives, put
\[
 u=(t/\alpha)^{1/m}>0,\qquad
 v=((1-t)/\beta)^{1/m}>0.
\]
Both roots are rational by construction, and the four ordinary points are
$[\pm u:\pm v:1]$, with independent signs. The point dictionary records
these exact rational coordinates. The program verifies their equations
using rational arithmetic. Each of the three, respectively fifteen,
class pairs satisfies a criterion; no two listed values give the same
pair. Thus their four point sections are complete for those pairs.

The three bits of \texttt{cusp\_flags} correspond to $0,1,\infty$.
They equal one precisely for
\[
 \beta=1,\qquad \alpha=1,\qquad \beta=-\alpha,
\]
respectively. Their two rational cusps are
$[0:\pm1:1]$, $[\pm1:0:1]$, and $[1:\pm1:0]$.
Each of these individual cusps has its complete Kummer packet. On the
tripod, the packet congruences are respectively
$[t]_m=[\alpha]_m$, $[1-t]_m=[\beta]_m$, and
$[1/t]_m=[\alpha^{-1}]_m$, as in A.23. The finite table counts curves
and packet centres, never the infinitely many sections in a packet.

The lossless decoding program is
\texttt{a\_main\_cm\_pairs.py}; for example, its
\texttt{sextic --decode a b} option returns the two class representatives,
all applicable targets, all ordinary coordinates, and all rational cusps
for identifiers $a,b$. Ordinary-point lists in rows with zero mask are
unclassified. A rational cusp flagged in such a row remains a known
cusp, but exhaustiveness is asserted only for nonzero masks. The exact
mask histograms, class and point dictionaries, and counts accompany
the full CSVs. A.36 supplies exhaustiveness of the section lists. The
generator iterates each possible identifier pair exactly once, so the
displayed counts are reproducible finite enumerations, not a sample.
\end{proof}

\begin{statement}{A.43}{The least conductor rank-one cubic twist}{Proved}
Among all cubic twists $E_d:y^2=x^3-432d^2$, with $d$ ranging over
positive cube-free integers, the least conductor of an analytic-rank-one
twist is $243$. It is attained by $d=9$, whose integral minimal model is
\[
 E_9^{\min}:\quad y^2+y=x^3-1.
\]
It has torsion-free Mordell--Weil group of rank one and contains the
non-torsion point $(1,0)$. Its first derivative has the certified enclosure
\[
 L'(E_9,1)\in
 [1.29019059036986358169134\ \mathbin{\pm}\ 6.10\cdot10^{-24}].
\]
This identifies the curve requested in task A.12; it does not determine
the profinite lifting sets posed there.
\end{statement}
\begin{proof}
Every rational cube class has a unique positive cube-free integral
representative, so this range is global. The displayed Weierstrass model
has discriminant
\[
 \Delta=-2^{12}3^9d^4,
\]
and $j=0$. The discriminant valuation modulo twelve is invariant under
an isomorphism over a local field. For every $p\mid d$, $p\ne3$, it is
$4$ or $8$ modulo twelve; for $p=2$ the extra initial twelve does not
change this conclusion. At $3$ its value is $9,13$, or $17$ modulo
twelve. Thus all these places have bad reduction.

Since $j$ is integral the curve has potentially good reduction. Choose
an auxiliary coefficient prime different from the residue
characteristic. Its inertia image on the rational Tate module is finite,
and its determinant on inertia is one. A nontrivial finite subgroup of
$\mathrm{SL}_2$ in characteristic zero has no invariant vector: an
element fixing a nonzero vector has both eigenvalues one and, being of
finite order, is the identity. Consequently the tame contribution to the
Artin conductor is two at each bad potentially good place. In particular,
\[
 f_3\geq2,\qquad f_p\geq2\quad(p\mid d,\ p\ne3).
\]
If $3\mid d$, wild inertia is nontrivial. Otherwise a finite tamely
ramified extension of ramification index $e$ prime to three would give
good reduction, by the inertia criterion. Discriminant scaling would
then require $12\mid e(9+4v_3(d))$, impossible since
$v_3(d)=1$ or $2$ and $3\nmid e$. The Swan term is therefore positive,
and $f_3\geq3$ in this case. This argument uses only the potential-good
reduction and inertia criteria, rather than a guessed complete
conductor formula; see the local representation framework of
\citet{AMSerre1972}.

If $N\leq243$, no prime $p\geq7$ can divide $d$, since then
$N\geq9p^2>243$. If five divides $d$ together with two, then
$N\geq9\cdot4\cdot25>243$; if five and three divide $d$, then
$N\geq27\cdot25>243$. Hence every possible $d$ belongs to the finite
list supported on two and three, together with $5,25$. Their exact
conductors and ranks, already certified in A.38, are
\[
\begin{array}{c|rrrrrrrrrrr}
d&1&2&3&4&5&6&9&12&18&25&36\\\hline
N&27&36&243&108&675&972&243&972&972&225&972\\
r_{\rm an}&0&0&0&0&0&1&1&1&0&0&0.
\end{array}
\]
Thus the only rank-one member of conductor at most $243$ is $d=9$.
The exact change of variables from its short model is
$x=36x'$, $y=216y'+108$, giving the stated minimal model. Its torsion
is trivial by A.39, so $(1,0)$ is non-torsion. A.37 certifies the first
derivative enclosure, and Rubin's CM rank-one theorem identifies the
algebraic rank with the analytic rank. No assertion that $(1,0)$ is a
chosen primitive generator is needed here.
\end{proof}
